% Transcription — fonds Alexandre Grothendieck, shelfmark 12, batch 4 (pages 61-80).
% Established from the facsimile archives/batches/12/batch-04.pdf (a mirror of
% https://grothendieck.umontpellier.fr/12.pdf), per the skill
% transcribe-grothendieck. The batch file carries no cover sheet: PDF page N
% is page 60+N of the folder. The Montpellier footer printed under each leaf
% confirms it (« 65 » under PDF page 5, « 70 » under PDF page 10), and so do
% the archivists' pencil numbers bottom-left, which agree with it on every
% leaf (61 on the first, 80 on the last).
%
% Pass: Opus 5.5 (claude-opus-5-5), 2026-09-23 — first pass, unchecked against the pages by a human.
%
% Transcribed: 61-64, 66, 68, 70, 72-79 (fifteen pages). Skipped: 65, 67, 69,
% 71, 80 — typescript of EGA (IV, § 7.5 on excellent rings and the
% conditions (i)-(viii), § 11.9, and a run of Propositions/Corollaires on
% F-regular elements and associated prime cycles), cancelled by long
% diagonal strokes; 71 is 65 showing through in mirror image. 80 carries two
% short handwritten notes in the margin of the typescript, about the
% typescript itself; like the EGA galley skipped in folder 13, it is left out.
% Page 79 is written in his hand on the back of that typescript (its text
% shows through) and is transcribed.
%
% Reading order. Page 76 ends « et le » and the sentence resumes at the head
% of 78; 77 is an insert called from the left margin of 78 by a matching
% sign (a stroke ending in a small circle). The \page{} order stays that of
% the leaves, 76-77-78, and a \note on 77 says so.
%
% What the batch is. The continuation of his numbered paragraph « 15)
% Cohomologie absolue », begun on page 59 (batch 3), then paragraph « 16 »,
% which he heads himself. No pagination of his and no date on any leaf.
%   61      end of 2°) (transitivity RΓ_X = RΓ_Y Rf and the Leray spectral
%           sequence); 3°) H^i(X,M) = 0 for i < 0, H^0(X,M) = Hom(Q_X(0),M);
%           N.B. on hyper-Ext, Ext^i(X;M,N) = H^i(X, RHom(M,N)), Ext^0 = Hom,
%           the hope that Ext^1 classifies extensions.
%   62-63   whether Ext^* is the Ext of the category of motives (doubts for
%           Ind-motives); spectral sequences H^p(X, H^q(M^·)) ⇒ H^*(X,M^·)
%           and H^p(X, Ext^q(M,N)) ⇒ Ext^*; 4°) the ℓ-adic realisation maps
%           H^i(X,M) ⊗_Q Q_ℓ → H^i(X, T_ℓ(M)), and Ext likewise, wanted to be
%           isomorphisms for X of finite type over Z, with H^*(X,M) of finite
%           type over Q; a struck box (the two cases X of finite type over Z
%           or over an algebraically closed field); over k algebraically
%           closed, H^i(k,M) = 0 for i > 0; an NB on explicit computation over
%           the closure of a finite field via the Tate conjectures.
%   63-64   Tate's theorem reformulated: a^i(X) ⊗ Q ≃ Γ_k(R f_*(Q_X(i))),
%           and a^i(X) ⊗ Q ≃ Im[H^{2i}(X, Q(i)) → H^{2i}(X, Q_ℓ(i))] (boxed);
%           over the closure of F_p, H^i(X, Q(j)) is 0 except in degrees 2j
%           (a^j(X) ⊗ Q, canonically) and 2j+1 (a^j(X) ⊗ Q, not canonically);
%           the image in ℓ-adic cohomology vanishes off degree 2i by weights.
%   66-68   « th. de dualité absolue »: M and M* dual under the dualising
%           motive Q(n)[2n+1] (X regular of dimension n), RΓ_X(M) and
%           R_!Γ_X(M*) dual; boxed: H^i(X,M) dual to Ext_!^{2n+1-i}(X;M,Q(n)),
%           H^i_!(X,M) dual to Ext^{2n+1-i}(X;M,Q(n)); via relative duality,
%           reduction to X = Spec Z and then to Spec F_p, i.e. the duality of
%           H^i(F_p,M) and H^{1-i}(F_p,M*) by cup-product and H^1(F_p,Q(0)) ≃ Q,
%           checked ℓ-adically as the triviality H^i(π,V) × H^{1-i}(π,V*).
%   70      X projective smooth geometrically connected over a finite field:
%           the Hochschild-Serre sequence, H^i(X,Q_ℓ(j)) = 0 off 2j, 2j+1,
%           a^j(X/k) ⊗ Q_ℓ in both via Tate (and semi-simplicity for 2j+1);
%           passage to the limit over the extensions k_n.
%   72      weight bookkeeping for a pure motive M of weight ρ over S smooth
%           of finite type over F_p: H^i(S,M_ℓ) = 0 unless
%           -2i ≤ ρ ≤ -i+1 (i ≥ 2), -2 ≤ ρ ≤ 1 (i = 1), ρ = 0 (i = 0), with
%           interval sketches in the margin.
%   73-79   « 16 Relations avec les pts rationnels (et la cohomologie) des
%           variétés abéliennes sur un schéma de type fini »: A_K over K of
%           finite type, spread out to an abelian scheme A over S regular of
%           finite type over Z, A_K(K) = A(S) finitely generated
%           (Mordell-Weil-Néron-Lang); the motive R^1f_*(Q_A(1)) ≃ T_mot(A);
%           the heuristic A(S) ⊗ Q ≃ H^1(S, T_mot(A)), H^0 = 0 by weights;
%           Kummer sequences; the boxed expectation H^1(S,A)(ℓ) finite for ℓ
%           prime to S (true for S = Spec F_q by Lang); over F_p, the diagram
%           (D) comparing S with S̄ via M = H^0(S̄,Ā)/H^0(F̄_p, A^{tr}), whose
%           last vertical arrow is bijective, so the conjecture becomes
%           H^1(S̄,Ā)(ℓ)^π finite (boxed), roughly equivalent (via Jacobians)
%           to H^2(S,G_m) → H^2(X,G_m) surjective mod finite groups for
%           X → S proper; NB on H^i(S,A)(ℓ) for i ≥ 3; the counterexample
%           S a curve, A = J × S, H^2(S,T_ℓ(A)) ≃ End(J) ⊗ Q_ℓ ≠ 0; page 79:
%           the conjecture means rank_{Q_ℓ} H^1(S̄,T_ℓ(Ā))^π = rank_Z A(K),
%           « essentiellement équivalente à la conjecture de
%           Birch-Swinnerton-Dyer ».
%
% Notation fixed for the batch, following batch 3: \mathbb{R}\Gamma_{X}, his
% double-struck R; H^i(X,M) absolute cohomology; \mathbb{E}\mathrm{xt} and
% \mathbb{H} for his double-struck hyper-Ext and hypercohomology;
% \underline{\mathrm{Ext}} his underlined sheaf Ext; T_ℓ the ℓ-adic
% realisation, T_{\mathrm{mot}}(A) the motive of A; a^i(X) his rounded a for
% the group of cycle classes; π = Gal(F̄_p/F_p) (= \hat{Z}); lower π for
% coinvariants, upper π for invariants; (ℓ) for the ℓ-primary part; _{ℓ^ν}A,
% _{ℓ^∞}A the kernels of ℓ^ν and their union; A^{\mathrm{tr}} as he writes it.
%
% The dating is the inventory's own for the folder, copied verbatim. Its
% group, [10-18] « Théorie des motifs », is dated
% [à partir de 1958-à partir de 1982].
\documentclass[11pt,a4paper]{article}
% Le préambule commun à toutes les transcriptions du fonds Grothendieck.
%
% Il tient en une idée : **l'incertitude se compose, elle ne se résout pas.**
% Une transcription qui lisse un mot douteux a détruit la seule chose pour
% laquelle on la faisait — un lecteur doit pouvoir savoir, à chaque endroit,
% ce qui était sur le feuillet et ce qui a été ajouté. Les six macros qui
% suivent sont donc l'appareil critique, et non de la décoration.
%
% Les mêmes commandes sont reconnues par `scripts/render.mjs`, qui produit la
% vue de lecture du site. Une macro ajoutée ici doit l'être aussi là-bas,
% faute de quoi le rendu échouera bruyamment — ce qui est voulu : un rendu
% silencieusement faux est pire qu'un rendu absent.

\ProvidesPackage{grothendieck}[2026/08/08 Transcriptions du fonds Grothendieck]

\usepackage{amsmath,amssymb,amsthm}
\usepackage{mathrsfs}
\usepackage{stmaryrd}
% Les diagrammes commutatifs sont partout dans ce fonds : c'est la langue
% même de la géométrie algébrique de Grothendieck, et une transcription qui
% les rendrait en prose ne transcrirait rien.
\usepackage{tikz-cd}
% Français par défaut — c'est la langue des feuillets — mais l'anglais est
% chargé aussi : la traduction et le résumé sont des documents anglais, et la
% typographie française (espaces avant les deux-points) y serait une faute.
% Ces documents basculent par \selectlanguage{english} en tête de corps.
\usepackage[english,french]{babel}
\usepackage{fontspec}
\usepackage{geometry}
\geometry{a4paper,margin=25mm,marginparwidth=28mm}
\usepackage{marginnote}
\usepackage{xcolor}
% ulem, not soul. `soul`'s \st and \ul are built on a character-by-character
% scan that XeLaTeX's font machinery breaks: the document fails with an
% undefined control sequence rather than degrading. ulem does the same two jobs
% and is Unicode-engine safe. `normalem` stops it hijacking \emph, which the
% transcriptions use for Grothendieck's own emphasis.
\usepackage[normalem]{ulem}
\usepackage{hyperref}
\hypersetup{colorlinks=true,linkcolor=black,urlcolor=black,pdfborder={0 0 0}}

% --- Métadonnées du lot -----------------------------------------------------
% Reprises telles quelles par le script de rendu, qui les affiche en tête de
% la vue de lecture. Elles fixent ce que le fichier prétend couvrir : sans
% elles, une transcription flotte sans point d'attache dans un fonds de seize
% mille feuillets.

\newcommand{\folder}[1]{\gdef\grothFolder{#1}}
\newcommand{\batch}[1]{\gdef\grothBatch{#1}}
\newcommand{\leaves}[2]{\gdef\grothFirst{#1}\gdef\grothLast{#2}}
\newcommand{\foldertitle}[1]{\gdef\grothFoldertitle{#1}}
\gdef\grothFolder{?}\gdef\grothBatch{?}\gdef\grothFirst{?}\gdef\grothLast{?}\gdef\grothFoldertitle{}

% --- Le résumé --------------------------------------------------------------
%
% Il ouvre la lecture modernisée, sous le titre « Résumé », et il est composé
% en petit. Ce n'est pas un ornement : c'est le seul passage du document qui
% ne soit pas des mathématiques, et un lecteur doit pouvoir le reconnaître
% avant de l'avoir lu — puis le sauter s'il connaît déjà le sujet, ou s'y
% arrêter s'il ne le connaît pas. Le filet qui le suit marque la reprise du
% corps.

\newenvironment{resume}
  {\small\setlength{\parskip}{0.45em}}
  {\par\medskip\hrule\medskip}

% --- Mots-clés ---------------------------------------------------------------
%
% En anglais, en clôture du résumé de la lecture modernisée : le vocabulaire
% moderne sous lequel chercher ce que les feuillets construisent. Le manifeste
% du site (`npm run manifest`) les extrait pour étiqueter la cote entière —
% cette ligne est la source unique des tags, il n'y a pas de fichier à côté.
\newcommand{\keywords}[1]{\par\smallskip\noindent
  {\footnotesize\textsc{Keywords} --- \textit{#1}}\par}

% --- L'appareil critique ----------------------------------------------------

% Le numéro de feuillet, en marge. C'est l'ancre qui permet de régler un
% désaccord sur une formule en regardant *un* feuillet plutôt qu'une chemise
% de six cents. Le site s'en sert aussi pour tourner les pages du fac-similé
% quand on fait défiler la transcription.
\newcommand{\leaf}[1]{%
  \marginnote{\footnotesize\color{gray!70}\textsf{#1}}%
  \label{leaf:#1}\ignorespaces}

% Illisible : on ne devine pas. Un mot inventé qui se lit comme les autres est
% le pire résultat possible.
\newcommand{\ill}{\textcolor{red!60!black}{[\,\ldots\,]}}

% Lecture douteuse : le mot est donné, et signalé comme incertain.
\newcommand{\uncertain}[1]{\uline{#1}}

% Ajout de l'éditeur — une lettre manquante, un mot sous-entendu.
\newcommand{\add}[1]{\textcolor{blue!50!black}{[#1]}}

% Biffé par Grothendieck lui-même. Ce qu'il a rayé fait partie du document :
% une rature dit souvent où la pensée a changé de direction.
\newcommand{\struck}[1]{\sout{#1}}

% Note du transcripteur, sur l'état matériel du feuillet.
\newcommand{\note}[1]{\par\smallskip{\small\color{gray!60!black}\textit{[#1]}}\par\smallskip}

% Note marginale de Grothendieck, distincte du corps du texte.
\newcommand{\marginal}[1]{\par\smallskip\noindent\hspace{1em}%
  {\small\color{gray!60!black}#1}\par\smallskip}

% --- Titre ------------------------------------------------------------------

\AtBeginDocument{%
  \begin{center}
    {\large\bfseries Fonds Alexandre Grothendieck --- cote n\textsuperscript{o}~\grothFolder}\\[2pt]
    {\itshape\grothFoldertitle}\\[4pt]
    {\small feuillets \grothFirst--\grothLast\ (lot \grothBatch)}\\[2pt]
    {\footnotesize\color{gray!60!black}Transcription établie d'après le fac-similé numérisé
     par l'Université de Montpellier.\\
     Les lectures incertaines sont soulignées, les illisibles notées [\,\ldots\,],
     les ajouts entre crochets.}
  \end{center}
  \vspace{1em}}

\endinput


\folder{12}
\batch{4}
\pages{61}{80}
\dating{1964-[vers 1971]}
\watermark{Édition de démonstration}
\foldertitle{Généralités (sous-groupes de Galois motiviques) (1964 ?) : notes manuscrites (s.d.), tapuscrit (s.d.).}

\begin{document}

\page{61}
\note{la page continue le 2°) de la page 59 (lot précédent) : « si $f : X \to Y$ est » …}

type fini, on a
\[
  \mathbb{R}\Gamma_{X} = \mathbb{R}\Gamma_{Y}\, \mathbb{R}f_{*}
\]
d'où des suites spectrales dans $(\mathrm{Ab})$,
\[
  H^{*}(X, M) \Longleftarrow H^{p}(Y, R^{q}f_{*}(M))
\]

3°) Bien entendu, \uncertain{on a} \ill{} $H^{i}(X, M) = 0$ si $i < 0$, et
$H^{0}(X, M)$ \uncertain{est exact} \ill{} \uncertain{en} $M$. De plus,
\uncertain{dans} \uncertain{fonctoriel}
\[
  H^{0}(X, M) = \operatorname{Hom}(\mathbb{Q}_{X}(0), M)
\]
\note{dans $\mathbb{Q}_{X}(0)$, la lettre $\mathbb{Q}$ est surchargée (un $R$ repris en $\mathbb{Q}$ ?) et entourée.}

[N.B. Rien n'empêche d'introduire des hyper-\uncertain{ext}
\[
  \mathbb{E}\mathrm{xt}^{i}(X ; M, N) = \mathbb{H}^{i}(X, \mathbb{R}\underline{\mathrm{Hom}}(M, N)) .
\]
Cela redonne bien, par le n° 7,
\[
  \operatorname{Hom}(M, N) \simeq \mathbb{E}\mathrm{xt}^{0}(X, M, N)
\]
et on espère que $\mathrm{Ext}^{1}(X ; M, N)$ est bien l'ensemble des
\emph{classes d'extension} de $M$ par $N$. [On aimerait pouvoir dire que

\page{62}
$\mathbb{E}\mathrm{xt}^{*}(X ; M, N)$ est le \struck{\ill{}} Ext dans la
catégorie des motifs, i.e.\ aux foncteurs dérivés ?.
Pourtant, \uncertain{doutes} pour les Ext$^{i}$ ($i > 0$) \ill{}
\uncertain{pour être} dans la catégorie des Ind-motifs ???]
\uncertain{Possibles} …
\note{les deux lignes « doutes … Ind-motifs ??? » sont serrées l'une sur l'autre, la seconde en interligne ; le mot lu « Possibles » est traversé d'un grand trait vertical.}

Notons la suite spectrale
\[
  \mathbb{H}^{*}(X, M^{\bullet}) \Longleftarrow H^{p}(X, H^{q}(M^{\bullet}))
\]
d'où en particulier
\[
  \mathbb{E}\mathrm{xt}^{*}(X, M, N) \Longleftarrow H^{p}(X, \underline{\mathrm{Ext}}^{q}(M, N))
\]

4°) On veut \struck{que \ill{} des conditions standard}
\uncertain{bien le} \uncertain{$\ell$-foncteur} \struck{sur $X$}
\note{la ligne est biffée d'un long trait de « que » à « standard » ; au-dessus, en interligne, trois mots, dont on croit lire « dans », « bien le » et « $\ell$-foncteurs » ; la lecture de l'ensemble n'est pas sûre.}
\[
  H^{i}(X, M) \otimes_{\mathbb{Q}} \mathbb{Q}_{\ell} \longrightarrow H^{i}(X, T_{\ell}(M))
\]
\note{le $\otimes_{\mathbb{Q}}$ est repassé et à demi barré.}
d'où aussi
\[
  \mathrm{Ext}^{i}(X ; M, N) \otimes_{\mathbb{Q}} \mathbb{Q}_{\ell} \longrightarrow \mathbb{E}\mathrm{xt}^{i}(X ; T_{\ell}(M), T_{\ell}(N))
\]
qui, sous l'hypothèse que $X$ est de type fini sur $\mathbb{Z}$,
\struck{dans les \uncertain{conditions}, d'excellence} \struck{\ill{}}
(en invoquant le fait que les seconds membres sont \uncertain{de type fini}
sur $\mathbb{Q}_{\ell}$)
\note{un double trait vertical, dans la marge, court le long de ces lignes et des trois premières de la page suivante.}

\page{63}
soient des isomorphismes, avec le complément que $H^{*}(X, M)$ est un
\emph{$\mathbb{Q}$-Module de type fini}.
\struck{Ceci doit être \uncertain{vérifié} en particulier dans les deux cas
suivants : a) $X$ de type fini sur $\mathbb{Z}$ ; b) $X$ de type fini sur un
corps alg.\ clos.}
\note{ces deux lignes et la phrase qui les introduit sont encadrées, et le cadre est barré de grandes croix.}

En particulier, pour un corps alg.\ clos, on aurait
\[
  H^{i}(k, M) = 0 \quad \text{pour } i > 0, \quad \forall \text{ motif } M .
\]

\emph{NB} Si $X$ est de type fini (ou projectif lisse) sur un corps alg.\
clos $K$, \ill{} \ill{} \ill{} \ill{} telles relations précises, comme on
voit par calcul explicite (utilisant les conjectures de Tate) pour $K$
clôture algébrique d'un corps fini …

Le th.\ de Tate peut se reformuler de façon précise en énonçant :
\[
  a^{i}(X) \otimes_{\mathbb{Z}} \mathbb{Q} \xrightarrow[\;\sim\;]{\text{ét.\ pr.\ } \mathbb{Q}_{\ell}} \Gamma_{k}\bigl(R^{i}f_{*}(\mathbb{Q}_{X}(i))\bigr)
\]
\note{l'étiquette de la flèche, en exposant, se lit « ét.\ pr.\ $\mathbb{Q}_{\ell}$ » sans certitude ; l'exposant de $R$ est lu $i$, alors que la formule encadrée de la page suivante porte $H^{2i}$.}

\page{64}
$X$ étant projectif et lisse sur $k$ alg.\ clos, \uncertain{simplifiant}
\uncertain{les notations} ($f : X \to \operatorname{Spec} k$).
Une autre façon de l'écrire :
\[
  \boxed{\; a^{i}(X) \otimes_{\mathbb{Z}} \mathbb{Q} \simeq \operatorname{Im}\bigl[ H^{2i}(X, \mathbb{Q}(i)) \longrightarrow H^{2i}(X, \mathbb{Q}_{\ell}(i)) \bigr] \;}
\]
\note{le $\otimes_{\mathbb{Z}} \mathbb{Q}$ est surchargé et en partie barré ; les deux exposants $2i$ sont repris sur un premier $i$.}

\emph{NB.} Si $k$ est la clôture algébrique de $\mathbb{F}_{p}$, on trouve
\ill{} \ill{} pour $X$ projectif lisse sur $k$
\[
  H^{i}(X, \mathbb{Q}(j)) =
  \begin{cases}
    0 & \text{si } i \neq 2j, 2j+1 \\
    a^{j}(X) \otimes_{\mathbb{Z}} \mathbb{Q} & \text{si } i = 2j \quad \text{(canonique)} \\
    a^{j}(X) \otimes_{\mathbb{Z}} \mathbb{Q} & \text{si } i = 2j+1 \quad \text{(\emph{pas} canonique)}
  \end{cases}
\]
\marginal{pas $\mathbb{Q}_{\ell}(i)$ !}
\note{ajout au-dessus du $\mathbb{Q}(j)$ du premier membre.}

De plus, pour $i \neq 2j$, l'image de $H^{i}(X, \mathbb{Q}(j))$ dans
$H^{i}(X, \mathbb{Q}_{\ell}(j))$ est nulle.

C'est un \emph{fait général} : si $k$ alg.\ clos, \uncertain{dans}
($X$ projectif lisse sur $k$) l'image de
\[
  H^{j}(X, \mathbb{Q}(i)) \longrightarrow H^{j}(X, \mathbb{Q}_{\ell}(i))
\]
est \emph{nulle} pour $j \neq 2i$ (question de poids), pour $j = 2i$
\uncertain{c'est} les cycles algébriques \add{\uncertain{à coeff.}}
rationnels.
\note{les indices changent de rôle d'une formule à l'autre : $H^{i}(X, \mathbb{Q}(j))$ ci-dessus, $H^{j}(X, \mathbb{Q}(i))$ ici ; transcrits comme écrits. Le dernier ajout, en interligne au-dessus d'un mot biffé, se lit mal.}

\page{66}
Notons le th.\ de \uncertain{dualité} \uncertain{absolue} :

Soient $M$ et $M^{*}$ deux \add{\uncertain{(constructibles)}} motifs sur $X$
(de type fini sur $\mathbb{Z}$), qui sont « duaux » par le
\uncertain{couplage} de motifs dualisants \struck{\ill{}} canonique ---
[NB. Si $X$ \uncertain{est régulier} de dim.\ $n$, \uncertain{dans}
$\mathbb{Q}(n)[2n+1]$, c'est le dualisant \emph{à valeurs} dans
\struck{$\mathbb{Q}(n)[2n+1]$}]
\marginal{shift}
\note{« shift » est écrit deux fois dans la marge de droite, avec des flèches vers le crochet $[2n+1]$.}
\note{un mot entouré, en interligne dans le crochet, se lit « \uncertain{foncteur} » ; sa place dans la phrase n'est pas sûre. Le second $\mathbb{Q}(n)[2n+1]$, en fin de ligne, est biffé et surchargé.}

alors $\mathbb{R}\Gamma_{X}(M)$ et $\mathbb{R}_{!}\Gamma_{X}(M^{*})$ sont
duaux.
\marginal{\uncertain{supposer qu'on sait que} $\mathbb{Z}$}
\note{ligne ajoutée entre guillemets sous la formule.}
Ici [comme les coefficients sont ess.\ $\mathbb{Q}$, donc il n'y a pas de
Tor] ceci revient aussi à la formule (cas régulier)
\[
  \boxed{\;
  \begin{aligned}
    &H^{i}(X, M) \ \text{dual de}\ \mathrm{Ext}_{!}^{2n+1-i}(X ; M, \mathbb{Q}(n)) \\
    &H^{i}_{!}(X, M) \ \text{dual de}\ \mathrm{Ext}^{2n+1-i}(X ; M, \mathbb{Q}(n))
  \end{aligned}
  \;}
\]
Si $X$ propre sur $\operatorname{Spec} \mathbb{Z}$, inutile de mettre des
$!$.

Moyennant le th.\ de dualité \emph{relative} pour motifs, cet énoncé revient
au

\page{68}
cas choisi pour $X = \operatorname{Spec} \mathbb{Z}$, et lui-même
\uncertain{le} cas géométrique, pour $X = \operatorname{Spec} \mathbb{Z}/p\mathbb{Z}$.
Dans ce dernier cas, on doit vérifier simplement la dualité entre
\[
  H^{i}(\mathbb{F}_{p}, M) \quad \text{et} \quad H^{1-i}(\mathbb{F}_{p}, M^{*})
\]
\[
  \bigl( M^{*} = \underline{\mathrm{Hom}}(M, \mathbb{Q}(0)) \bigr)
\]
à l'aide du « cup-produit » et de l'isom.
\[
  H^{1}(\mathbb{F}_{p}, \mathbb{Q}(0)) \simeq \mathbb{Q}
\]
et il suffit de faire la vérification $\ell$-adiquement, en utilisant le
foncteur $T_{\ell}$.
\marginal{à développer \uncertain{au sujet} \uncertain{des} cup-produits à coeff.\ « \uncertain{abstraits} » …}
\note{note écrite en oblique dans la marge de gauche, à hauteur de « cup-produit ».}
\note{« cup-produit » est souligné d'un trait ondulé.}

Cela revient à la vérification suivante : [pour un \uncertain{vectoriel}
\ill{} de dim.\ finie sur $\mathbb{Q}_{\ell}$ sur lequel
$\pi = \widehat{\mathbb{Z}}$ opère]
\[
  H^{i}(\pi, V) \quad \text{et} \quad H^{1-i}(\pi, V^{*}) \quad \text{en dualité}
\]
$(V^{*} = \operatorname{Hom}_{\mathbb{Q}_{\ell}}(V, \mathbb{Q}_{\ell}))$ : c'est
une trivialité.

\page{70}
$X$ projective lisse \add{géom.\ connexe} sur $k$ corps fini à $q$ éléments.
Par passage à la limite sur les \struck{quotients} \add{coeff.} finis, on a
\[
  \longrightarrow H^{i-1}(\overline{X}, \mathbb{Q}_{\ell}(j))_{\pi} \longrightarrow H^{i}(X, \mathbb{Q}_{\ell}(j)) \longrightarrow H^{i}(\overline{X}, \mathbb{Q}_{\ell}(j))^{\pi} \longrightarrow 0
\]
\note{l'exposant du premier terme est surchargé ; lu $i-1$ d'après la suite.}
d'où le résultat
\[
  \begin{aligned}
    &\text{si } i \neq 2j, 2j+1, && H^{i}(X, \mathbb{Q}_{\ell}(j)) = 0 \\
    &\text{si } i = 2j, && H^{2j}(X, \mathbb{Q}_{\ell}(j)) = H^{2j}(\overline{X}, \mathbb{Q}_{\ell}(j))^{\pi} \underset{\text{conj.\ de Tate}}{\simeq} a^{j}(X/k) \otimes_{\mathbb{Z}} \mathbb{Q}_{\ell} \\
    &\text{si } i = 2j+1, && H^{2j+1}(X, \mathbb{Q}_{\ell}(j)) = H^{2j}(\overline{X}, \mathbb{Q}_{\ell}(j))_{\pi} \underset{\text{Tate} + \text{action semi-simple}}{\simeq} a^{j}(X/k) \otimes_{\mathbb{Z}} \mathbb{Q}_{\ell}
  \end{aligned}
\]
\note{les deux isomorphismes sont notés « $\wr\!\wr$ » sous les termes, avec leur justification en petit à côté ; dans la seconde, « action » est une lecture incertaine. Dans les deux $\otimes$ l'indice est surchargé.}

Soit maintenant $k_{n}$ l'extension de degré $n$ de $k$ dans $\overline{k}$,
et étudions
\[
  \varinjlim H^{i}(X_{k_{n}}, \mathbb{Q}_{\ell}(j))
\]
on trouve $0$ pour $i \neq 2j, 2j+1$ ; pour $i = 2j$, on trouve
$H^{2j}(\overline{X}, \mathbb{Q}_{\ell}(j))^{\mathfrak{J}} = a^{j}(\overline{X}) \otimes_{\mathbb{Z}} \mathbb{Q}_{\ell}$,
enfin pour $i = 2j+1$, on trouve également \add{\uncertain{en vertu du motif}}
$a^{j}(\overline{X}) \otimes_{\mathbb{Z}} \mathbb{Q}_{\ell}$.
\note{l'exposant de $H^{2j}(\overline{X}, \mathbb{Q}_{\ell}(j))$ est une lettre isolée, lue $\mathfrak{J}$ sans certitude.}
Donc on a
\[
  \varinjlim H^{i}(X_{k_{n}}, \mathbb{Q}(j)) =
  \begin{cases}
    0 & \text{si } i \neq 2j, 2j+1 \\
    a^{j}(\overline{X}) \otimes_{\mathbb{Z}} \mathbb{Q} & \text{si } i = 2j \\
    \struck{\ill{}} & \text{si } i = 2j+1, \ \text{mais pas canoniquement isomorphe à} \ldots
  \end{cases}
\]
et, sous la limite, $= H^{i}(\overline{X}, \mathbb{Q}(j))$.
\note{dans la troisième ligne, la valeur est biffée et surchargée (lue « $a^{j}(\overline{X})\otimes$ … » sous les ratures) ; la fin, en interligne, est coupée par le bord du feuillet et se termine sur « $X/\overline{\mathbb{F}}$ ! » (lecture incertaine) ; « $2j+1$ » est biffé à la fin de la ligne $i = 2j$.}
\marginal{N.B. \uncertain{Remplaçons} $k$ par \ill{} \ill{} \ill{} l'isom.\ obtenu par lim …}
\note{note serrée en bas à gauche du feuillet, en face de la dernière formule.}

\page{72}
$S$ \add{lisse} de type fini sur $\mathbb{F}_{p}$,
$M$ motif \ill{} de poids $\rho$
\[
  ?\ \ H^{*}(S, M_{\ell}) \overset{?}{\Longleftarrow} H^{p}(\mathbb{F}_{p}, R^{q}f_{*}(M_{\ell}))
\]
\[
  0 \longrightarrow \underbrace{H^{i-1}(\overline{S}, \overline{M}_{\ell})_{\pi}}_{\substack{\text{poids entre} \\ \rho + i - 1 \text{ et } \rho + \mathrm{Inf}(2i-2, 2n)}} \longrightarrow H^{i}(S, M_{\ell}) \longrightarrow \underbrace{H^{i}(\overline{S}, \overline{M}_{\ell})^{\pi}}_{\substack{\text{poids entre} \\ \rho + i \text{ et } \rho + \mathrm{Inf}(2i, 2n)}} \longrightarrow 0
\]
\note{l'exposant du premier terme est surchargé, lu $i-1$ ; le premier « $\mathrm{Inf}(2i-2, 2n)$ » se lit mal et est restitué d'après la ligne qui suit.}

Donc $H^{i}(S, M_{\ell}) = 0$ \struck{si on a} sauf si on a
\[
  \rho + i - 1 \leq 0 \leq \rho + 2i - 2 \quad \text{ou} \quad \rho + i \leq 0 \leq \rho + 2i
\]
i.e.
\[
  2 - 2i \leq \rho \leq 1 - i \quad \text{ou} \quad -2i \leq \rho \leq -i
\]
i.e.
\[
  \begin{aligned}
    -2i \leq \rho &\leq -i + 1 && (\text{si } i \geq 2) \\
    -2 \leq \rho &\leq 1 && (\text{si } i = 1) \\
    \rho &= 0 && (\text{si } i = 0)
  \end{aligned}
\]
\struck{Pour $i = 1$, cela donne} $\rho$
\note{dans « $2 - 2i \leq \rho \leq 1 - i$ », le $1 - i$ est récrit sur un premier terme biffé. Dans la marge de gauche, des croquis de segments : $[-2i+2, -i+1]$ et $[-2i, -i]$ superposés, puis un segment gradué $-2, -1, 0, 1$ ; un troisième, biffé, entre les deux.}

\section{16 Relations avec les points rationnels (et la cohomologie) des variétés abéliennes sur un schéma de type fini}
\note{titre de sa main en tête du feuillet 73, précédé du numéro « 16 » souligné ; « et la cohomologie » est ajouté en interligne, et « schéma de type fini … » remplace « corps de fonctions », biffé.}

\page{73}
Soit $A_{K}$ VA sur un corps $K$ de \emph{type fini} pour commencer.
Alors $K$ est le corps des fonctions d'un schéma $S$ irréductible et
régulier, de type fini sur $\mathbb{Z}$, \add{tel que} $A_{K}$ provient d'un
schéma abélien sur $S$. De plus, on a
\[
  A_{K}(K) = A(S)
\]
et on sait que c'est un \emph{groupe de type fini} par
Mordell-Weil-Néron-Lang.

Soit $f : A \to S$ la projection, et considérant le motif
\[
  \mathbb{R}^{1}f_{*}(\mathbb{Q}_{A}(1)) \simeq T_{\mathrm{mot}}(A)
\]
[\uncertain{où} $A^{*}$ \uncertain{est le} dual de $A$] \ill{} sur $S$
\struck{de} [qui est un motif] pur de poids $1$, qui redonne $A_{K}$ à
isogénie près \struck{\ill{}} (sur $K$ ?) comme conséquence des conjectures
de Tate.
\note{dans la formule, un indice $A$ sous le $\mathbb{Q}$ et un exposant $*$ sur le $f$ ; le second membre, « $T_{\mathrm{mot}}(A)$ », est récrit au-dessus d'un premier terme biffé. Le mot qui suit le premier crochet est le même que celui de la page 72 (« motif \ill{} de poids $\rho$ ») et reste illisible.}

\page{74}
On s'attend donc à pouvoir écrire $A(S) \otimes_{\mathbb{Z}} \mathbb{Q}$ en
termes de ce motif $M$.
Or en effet, des raisons heuristiques convaincantes pour nous suggèrent
\[
  A(S) \otimes_{\mathbb{Z}} \mathbb{Q} \simeq H^{1}(S, T_{\mathrm{mot}}(A))
\]
[NB. On doit avoir $H^{0}(S, T_{\mathrm{mot}}(A)) = 0$, à cause des
\struck{Mordell-Weil} considérations de poids …].
\struck{Une autre façon de} \uncertain{Considérons} \ill{} \ill{}
(\uncertain{aux} \ill{} \ill{} \ill{} $\times$)
\[
  0 \longrightarrow {}_{\ell^{\nu}}A \longrightarrow A \longrightarrow A \longrightarrow 0
\]
d'où
\[
  \longrightarrow H^{0}(S, A)_{\ell^{\nu}} \longrightarrow H^{1}(S, {}_{\ell^{\nu}}A) \longrightarrow H^{1}(S, A)_{\ell^{\nu}} \longrightarrow 0
\]
\note{les deux indices $\ell^{\nu}$ hors des parenthèses sont placés en bas à droite ; celui du troisième terme est récrit.}
et par passage à la limite et $\otimes_{\mathbb{Z}_{\ell}} \mathbb{Q}_{\ell}$
\[
  0 \longrightarrow \underbrace{H^{0}(S, A)}_{\text{de type fini sur } \mathbb{Z}} \otimes_{\mathbb{Z}} \mathbb{Q}_{\ell} \longrightarrow \underset{\substack{\wr\wr \\ H^{1}(S, T_{\mathrm{mot}}(A)) \otimes_{\mathbb{Q}} \mathbb{Q}_{\ell}}}{H^{1}(S, T_{\ell}(A))_{\mathbb{Q}_{\ell}}} \longrightarrow T_{\ell}(H^{1}(S, A))_{\mathbb{Q}_{\ell}} \longrightarrow 0
\]
\note{plusieurs indices de cette suite sont surchargés ($\otimes_{\mathbb{Z}} \mathbb{Q}_{\ell}$ récrit sur $\otimes \mathbb{Q}$, le $\ell$ de $T_{\ell}$ repassé).}

et le résultat escompté ici équivaut vraisemblablement à celui-ci :
\[
  \boxed{\; H^{1}(S, A)(\ell) \ \underline{\text{fini}} \ \text{si } \ell \text{ “premier à } S \text{”} \quad (S \text{ de type fini sur } \operatorname{Spec} \mathbb{Z}) \;}
\]
p.ex.\ si $S = \operatorname{Spec} \mathbb{F}_{q}$, alors c'est vrai par Lang !

\page{75}
Si $S$ est de type fini sur $\mathbb{F}_{p}$, on a (travaillant avec des
\emph{faisceaux étales}) si $f : S \to \operatorname{Spec} \mathbb{F}_{p}$
\[
  \longrightarrow H^{1}(\mathbb{F}_{p}, R^{0}f_{*}A) \longrightarrow H^{1}(S, A) \longrightarrow H^{0}(\mathbb{F}_{p}, \mathbb{R}^{1}f_{*}(A)) \longrightarrow \struck{\ill{}}
\]
\[
  H^{0}(\mathbb{F}_{p}, \mathbb{R}^{1}f_{*}(A)) = H^{1}(\overline{S}, \overline{A})^{\pi}
\]
\note{le dernier terme de la suite est un gribouillis biffé, où l'on devine un $\mathbb{F}_{p}$ ; l'égalité de la seconde ligne est notée par « $=$ » vertical sous le troisième terme.}

Or on a, toujours par « Kummer »
\[
  \longrightarrow \underset{\substack{\wr\wr \\ M \otimes \mathbb{Q}_{\ell}/\mathbb{Z}_{\ell}}}{H^{0}(\overline{S}, \overline{A}) \otimes \mathbb{Q}_{\ell}/\mathbb{Z}_{\ell}} \longrightarrow H^{1}(\overline{S}, {}_{\ell^{\infty}}\overline{A}) \longrightarrow H^{1}(\overline{S}, \overline{A})(\ell) \longrightarrow 0
\]
\note{« $M \otimes \mathbb{Q}_{\ell}/\mathbb{Z}_{\ell}$ » est écrit au-dessus du premier terme, et non au-dessous.}
d'où le fait que $H^{1}(\overline{S}, \overline{A})(\ell)$ est de poids
$1 - 1 = 0$ \add{si p.ex.\ $S$ est projectif lisse sur $\mathbb{F}_{p}$} --- à
priori rien ne s'oppose \uncertain{à ce qu'il y ait} un \uncertain{gros}
module d'invariants. D'autre part, on a
\[
  0 \longrightarrow \underbrace{A^{\mathrm{tr}}}_{\ill{}\ \text{de la V.A. } A \text{ sur } S} \longrightarrow \mathbb{R}^{0}f_{*}(A) \longrightarrow \underbrace{M}_{\substack{\pi\text{-module,} \\ \text{de type fini sur } \mathbb{Z}}} \longrightarrow 0
\]
d'où
\[
  \underset{\substack{\| \\ 0 \\ \text{par Lang}}}{H^{1}(\mathbb{F}_{p}, A^{\mathrm{tr}})} \longrightarrow H^{1}(\mathbb{F}_{p}, \mathbb{R}^{0}f_{*}(M)) \longrightarrow \underset{\substack{\| \\ (M \otimes \mathbb{Q}/\mathbb{Z})^{\pi}}}{H^{1}(\mathbb{F}_{p}, M)} \longrightarrow 0
\]
\note{le deuxième terme porte « $f_{*}(M)$ » là où l'on attend « $f_{*}(A)$ » ; transcrit comme écrit. Le dernier, « $(M \otimes \mathbb{Q}/\mathbb{Z})^{\pi}$ », est surchargé.}

\page{76}
Si $S$ de type fini sur $\mathbb{F}_{p}$, on obtient le diagramme suivant

\begin{tikzcd}[column sep=small, row sep=small, nodes={font=\scriptsize}]
 & M^{\pi} \otimes \mathbb{Q}_{\ell}/\mathbb{Z}_{\ell} \arrow[d, no head, "\wr" description] & & & \\
0 \arrow[r] & H^{0}(S, A) \otimes \mathbb{Q}_{\ell}/\mathbb{Z}_{\ell} \arrow[r] \arrow[d] & H^{1}(S, {}_{\ell^{\infty}}A) \arrow[r] \arrow[d] & H^{1}(S, A)(\ell) \arrow[r] \arrow[d] & 0 \\
0 \arrow[r] & (H^{0}(\overline{S}, \overline{A}) \otimes \mathbb{Q}_{\ell}/\mathbb{Z}_{\ell})^{\pi} \arrow[r] \arrow[d, no head, "\wr" description] & H^{1}(\overline{S}, {}_{\ell^{\infty}}\overline{A})^{\pi} \arrow[r] & H^{1}(\overline{S}, \overline{A})(\ell)^{\pi} \arrow[r] & 0 \\
 & (M \otimes \mathbb{Q}_{\ell}/\mathbb{Z}_{\ell})^{\pi} & & &
\end{tikzcd}

\note{dans la marge de gauche, le diagramme est nommé « (D) ».}

[où $M = H^{0}(\overline{S}, \overline{A}) / H^{0}(\overline{\mathbb{F}}_{p}, \overline{A}^{\mathrm{tr}})$,
$\pi$ est le groupe de Galois de $\overline{\mathbb{F}}_{p}$ sur
$\mathbb{F}_{p}$]

Les déterminations $H^{0}(S, A) \simeq M^{\pi}$ et
$H^{0}(\overline{S}, \overline{A}) \otimes \mathbb{Q}_{\ell}/\mathbb{Z}_{\ell} \simeq M \otimes \mathbb{Q}_{\ell}/\mathbb{Z}_{\ell}$
sont \uncertain{standard} ; on voit que la première flèche verticale est un
isomorphisme mod groupes finis. D'autre part, on trouve (suite spectrale de
Leray pour $f : S \to \operatorname{Spec} \mathbb{F}_{p}$ pour la
\emph{top.\ étale})
\[
  0 \longrightarrow \underset{\substack{\wr \\ \text{gr.\ fini}}}{H^{1}(\mathbb{F}_{p}, \mathbb{R}^{0}f_{*}(A))} \longrightarrow H^{1}(S, A) \longrightarrow \underset{\substack{\| \\ H^{1}(\overline{S}, \overline{A})^{\pi}}}{H^{0}(\mathbb{F}_{p}, \mathbb{R}^{1}f_{*}(A))} \longrightarrow 0
\]
\note{sous le premier terme, le signe lu « $\wr$ » peut être un simple trait de renvoi.}
comme on voit utilisant
\[
  0 \longrightarrow \widetilde{A}^{\mathrm{tr}} \longrightarrow \mathbb{R}^{0}f_{*}(A) \longrightarrow M \longrightarrow 0 ,
\]
et le

\page{77}
\note{ce feuillet ne continue pas le précédent, dont la phrase « et le » se poursuit en tête du feuillet 78. Il commence par un signe de renvoi (un trait terminé par un rond) qu'on retrouve, avec une flèche, dans la marge de gauche du feuillet 78, à hauteur de « $H^{2}(S, \mathbb{G}_{m})$ » : c'est un ajout appelé de là.}

$H^{2}(S, A)$ \struck{\ill{}} (donc le $T_{\ell}$ ou
$H^{2}(S, T_{\ell}(A))$) n'est pas nécessairement nul,
\add{comme on voit} \struck{\ill{}} si $S$ est une courbe lisse projective
géom.\ connexe sur $\mathbb{F}_{p}$ et $A = J \times S$, où $J$ est la
jacobienne de $S$. On a en tous cas
\[
  H^{2}(S, T_{\ell}(A)) \simeq H^{1}(\overline{S}, T_{\ell}(\overline{A}))_{\pi} = \bigl[ T_{\ell}(\overline{J}) \otimes T_{\ell}(\overline{J}) \bigr]_{\pi}
\]
\[
  \simeq \operatorname{End}(J) \otimes_{\mathbb{Z}} \mathbb{Q}_{\ell} \quad \bigl( \simeq (\text{classes de corr.\ divisorielles sur } S \times S) \otimes_{\mathbb{Z}} \mathbb{Q}_{\ell} \bigr)
\]
qui n'est pas nul ….
\note{« classes » est une lecture incertaine ; « divisorielles » est ajouté au-dessus de « corr. ». Avant le crochet, un premier second membre est biffé : « $= T_{\ell} J(\overline{\mathbb{F}}_{p}) \otimes$ » ; dans $H^{2}(S, T_{\ell}(A))$, le $T_{\ell}$ est récrit sur une première écriture ; avant la dernière parenthèse, un mot biffé.}

\page{78}
fait que $H^{1}(\mathbb{F}_{p}, A^{\mathrm{tr}}) = 0$ (Lang) et
$H^{1}(\mathbb{F}_{p}, M) = \text{fini}$ [trivial ; attention qu'il s'agit de
la coh.\ \emph{étale}, i.e.\ la coh.\ galoisienne de $\widehat{\mathbb{Z}}$ !],
donc la dernière flèche verticale du diagramme (D) est aussi bijective :
\[
  H^{1}(S, A)(\ell) \xrightarrow{\;\sim\;} H^{1}(\overline{S}, \overline{A})^{\pi}(\ell)
\]
\uncertain{sauf à noyau} \uncertain{fini} …
\note{dans le premier terme, le $(\ell)$ est récrit sur une première lettre.}

Donc la conjecture précédente signifie aussi
\[
  \boxed{\; H^{1}(\overline{S}, \overline{A})(\ell)^{\pi} = \text{groupe fini} \;}
\]

\struck{Je vois une façon de la prouver}
C'est \add{aussi} $\pm$ équiv.\ (via des jacobiennes) à ceci : $X \to S$
morphisme propre [\uncertain{lisse} de dim.\ relative $1$], $S$ régulier et
de type fini sur $\mathbb{Z}$. Prouver que
\[
  H^{2}(S, \mathbb{G}_{m}) \longrightarrow H^{2}(X, \mathbb{G}_{m})
\]
est surjectif mod groupes finis ??

\emph{NB} On devrait avoir $H^{i}(S, A)(\ell)$ fini aussi pour
$\underline{\underline{i \geq 3}}$, $S$ de type fini sur $\mathbb{Z}$, $\ell$
premier aux car.\ résiduelles
\add{NB} ($\text{si } i \geq 2$, $T_{\ell}(H^{i}(S, A)(\ell)) \simeq H^{i}(S, T_{\ell}(A))$)
\note{la dernière parenthèse est serrée en bas du feuillet et se lit mal ; le « $T_{\ell}$ » initial est ajouté sous la ligne.}
\marginal{donc les seuls \uncertain{cas} qui restent sont $H^{0}$ et $H^{2}$ pour $A$, \struck{\ill{}} $H^{1}$, $H^{2}$ pour $T_{\ell}(A)$}
\note{note écrite en oblique en bas à gauche, reliée par un trait au « $i \geq 3$ ».}

\page{79}
Notons que
\[
  H^{1}(S, T_{\ell}(A)) \simeq H^{1}(\overline{S}, T_{\ell}(\overline{A}))^{\pi}
\]
Donc notre conjecture signifie essentiellement
\[
  \operatorname{rang}_{\mathbb{Q}_{\ell}} \bigl( H^{1}(\overline{S}, T_{\ell}(\overline{A})) \bigr)^{\pi} \mathrel{\underline{\simeq}} \operatorname{rang}_{\mathbb{Z}} A(K)
\]
et sous cette forme, est essentiellement équivalente à la conjecture de
\struck{\ill{}} \emph{Birch-Swinnerton-Dyer}.
\note{le $T_{\ell}$ du premier membre est récrit sur un premier $T$. Le signe entre les deux rangs est un « $\simeq$ » souligné, qu'on peut lire « $=$ ».}

\end{document}
